\documentclass[12pt]{article}
\usepackage[utf8]{inputenc}
\usepackage[T1]{fontenc}
\usepackage[margin=1in]{geometry}
\usepackage{setspace}
\usepackage{xcolor}
\usepackage[colorlinks=true, urlcolor=blue!60!black, linkcolor=blue!60!black]{hyperref}
\setlength{\parindent}{2em}
\setlength{\parskip}{1.25em}
\renewcommand{\baselinestretch}{1.0}
\begin{document}
\begin{center}
{\LARGE \textit{United States v. Rahimi}, 602 U.S. 680 (2024)\par}
\vspace{0.8em}
{\large Case Brief --- Four-Page Edition\par}
\end{center}
\vspace{1em}

\section{Memory Jogger}

In an 8--1 decision clarifying \textit{Bruen}, the Supreme Court upheld 18 U.S.C.\ \S\ 922(g)(8) against a facial Second Amendment challenge, holding that an individual found by a court to pose a credible threat to the physical safety of another may be temporarily disarmed consistent with the principles underpinning the Nation's historical tradition of firearm regulation (\href{https://www.law.cornell.edu/supremecourt/text/22-915}{602 U.S. 680}).

\section{Detailed Case Facts}

In December 2019, Zackey Rahimi argued with his girlfriend, C. M., the mother of his child, in a Texas parking lot, dragged her into his car, fired a gun as she fled, and later threatened to shoot her if she told anyone. A Tarrant County court issued a two-year protective order, by consent and after notice and an opportunity to contest, finding that Rahimi had committed ``family violence'' and posed ``a credible threat'' to the ``physical safety'' of C. M.\ or their child; it suspended his gun license.

Rahimi repeatedly violated the order, threatened another woman with a gun, and was linked to five shootings in Arlington, including firing at a motorist and into the air at a burger restaurant. A search of his home turned up a pistol, a rifle, and a copy of the order.

\section{Procedural History}

Rahimi was indicted in the Northern District of Texas under \S\ 922(g)(8), which applies to orders issued after notice and a hearing that restrain threats against an intimate partner or child and either (i) find a ``credible threat to the physical safety'' of that person or (ii) prohibit the use of force. Bound by circuit precedent, the district court denied his facial challenge, and he pleaded guilty. After \textit{Bruen} (2022) replaced means-end scrutiny with a text-and-history test, the Fifth Circuit vacated the conviction for want of a historical analogue. The Supreme Court reversed on June 21, 2024.

\section{Judicial Votes}

Reversed and remanded, \textbf{8--1}.
\begin{itemize}
\item \textbf{Majority:} Chief Justice Roberts, joined by Justices Alito, Sotomayor, Kagan, Gorsuch, Kavanaugh, Barrett, and Jackson.
\item \textbf{Concurrences:} Sotomayor (joined by Kagan); Gorsuch; Kavanaugh; Barrett; Jackson.
\item \textbf{Dissent:} Justice Thomas.
\end{itemize}

\section{Holding}

An individual found by a court to pose a credible threat to the physical safety of another may be temporarily disarmed. Section 922(g)(8)(C)(i) is constitutional as applied to Rahimi and so survives his facial challenge under \textit{United States v. Salerno}. Under \textit{Bruen}, a modern law need not be a ``dead ringer'' or ``historical twin''; courts ask whether it fits the \textit{principles} underlying the tradition, examining ``why'' and ``how'' it burdens the right. The Government may not disarm people merely as not ``responsible.'' The Court reserved \S\ 922(g)(8)(C)(ii), categorical bans such as \S\ 922(g)(1), and whether 1791 or 1868 controls.

\section{Analysis of Opinions}

\textbf{Majority (Roberts).} The Court's precedents ``were not meant to suggest a law trapped in amber''; just as the right covers modern arms, it permits regulations ``relevantly similar'' to historical principles. Surety laws let magistrates, after notice and a hearing, require suspected abusers to post peace bonds, and ``going armed'' laws descended from the Statute of Northampton punished terrifying the public with weapons. Section 922(g)(8)(C)(i) shares their ``why'' (mitigating demonstrated threats) and ``how'' (temporary disarmament after an individualized judicial finding, a lesser burden than jail). The Fifth Circuit wrongly demanded a ``historical twin'' and inverted \textit{Salerno} by chasing hypothetical unconstitutional applications, rather than asking whether the law is valid in at least some applications.

\textbf{Concurrences.} \textit{Sotomayor} (with Kagan) still thought \textit{Bruen} wrong and warned that Thomas's approach would make it a ``one-way ratchet,'' noting that an abused woman is five times likelier to be murdered when her abuser has a gun. \textit{Gorsuch} defended text and history against a ``black box regime'' of balancing and listed open as-applied questions. \textit{Kavanaugh} defended text, pre- and post-ratification history, and precedent over tiers of scrutiny. \textit{Barrett} limited binding history to original meaning and held that ``[h]istorical regulations reveal a principle, not a mold.'' \textit{Jackson} found ``little method to \textit{Bruen}'s madness.''

\textbf{Thomas (dissenting).} ``[N]ot a single historical regulation justifies'' the statute. Surety laws imposed a lesser burden, a bond that left arms intact, and affray laws were criminal prosecutions for public terror, not domestic conduct; he accused the majority of ``mixing and matching'' one law's justification with another law's burden. English ``dangerous person'' laws such as the Militia Act of 1662 targeted rebellion, the very abuse the Second Amendment answered, and States should disarm violent offenders through criminal conviction, not civil orders lacking Sixth Amendment safeguards.

\section{Examples: Future Applications}

\begin{enumerate}
\item \textbf{Same-Side (Upheld): Red-Flag Order.} After a noticed hearing, a judge issues a one-year order disarming a respondent found by clear and convincing evidence to pose a credible threat to others. It mirrors \S\ 922(g)(8)(C)(i).
\item \textbf{Same-Side (Upheld): Anti-Stalking Order.} A court finds a stalker a credible threat to a neighbor. \textit{Rahimi}'s principle covers threats ``to another,'' and surety laws protected neighbors too.
\item \textbf{Opposite-Side (Doubtful): Nonviolent Felon.} A \S\ 922(g)(1) prosecution of someone with one twenty-year-old food-stamp fraud conviction imposes permanent, categorical disarmament without any judicial finding of threat.
\item \textbf{Opposite-Side (Doubtful): ``Irresponsible'' Citizens.} A ban on those with three speeding tickets or a civil default rests on the ``not responsible'' theory the Court rejected.
\item \textbf{Fence-Sitter: (C)(ii) and Suicide-Only Orders.} A boilerplate divorce order barring force without a threat finding, or an order based solely on risk of self-harm, tests the reserved questions and the level of generality.
\end{enumerate}

\section{Critique}

\textbf{Originalist and textualist views.} Commentators saw ``one, two \dots many originalisms'' among the conservatives. Strict originalists echo Thomas that the Court practiced ``\textit{Bruen}-lite,'' moving from concrete traditions to abstract principles, and the Kavanaugh--Barrett split leaves unclear whether tradition has independent authority.

\textbf{Progressive and doctrinal views.} Joseph Blocher, Eric Ruben, and Reva Siegel argue that \textit{Rahimi} averted arming abusers but left \textit{Bruen}'s level-of-generality problem unsolved, that tying gun safety to an era that excluded women and tolerated abuse is structurally biased, and that choosing analogues involves inescapable value judgments.

\section{Key Quotations}

\begin{quote}
``Taken together, the surety and going armed laws confirm what common sense suggests: When an individual poses a clear threat of physical violence to another, the threatening individual may be disarmed.''
\end{quote}
\noindent --- Chief Justice Roberts, Opinion of the Court, slip op.\ at 13.

\begin{quote}
``Rather, we conclude only this: An individual found by a court to pose a credible threat to the physical safety of another may be temporarily disarmed consistent with the Second Amendment.''
\end{quote}
\noindent --- Chief Justice Roberts, Opinion of the Court, slip op.\ at 17.

\begin{quote}
``Nevertheless, some courts have misunderstood the methodology of our recent Second Amendment cases. These precedents were not meant to suggest a law trapped in amber.''
\end{quote}
\noindent --- Chief Justice Roberts, Opinion of the Court, 602 U.S.\ at 691.

\end{document}
